\chapter{The \texttt{algo} Command}
\label{ch:cli}

All work with \algolang\ outside notebooks goes through one command, \code{algo}. On Windows it is \file{algo.bat}; on Linux and macOS it is the shell script \file{algo}. Both use the private \file{.venv} created by the installer when it exists.

\begin{terminal}
$ algo --help
usage: algo [-h] [--version] command ...

AlgoLang 0.2.0 — every unit of work is an agent

positional arguments:
  command
    check     check a program for errors
    build     compile to Python
    run       compile and run flow main
    new       create a starter program
    doctor    check that the environment is ready
    notebook  open the AlgoLang notebook in Jupyter (installs Jupyter if needed)
\end{terminal}

\section{Commands}

\begin{center}
\small
\begin{tabularx}{\linewidth}{@{}L{5.2cm}Y@{}}
\toprule
\thead{Command} & \thead{What it does}\\
\goldrule
\code{algo check FILE.algo} & Parses and checks the program; reports all errors; writes nothing\\
\code{algo build FILE.algo [-o OUT.py]} & Compiles to Python; default output \file{build/FILE.py} next to the source\\
\code{algo run FILE.algo [options]} & Compiles, then runs \code{flow main}\\
\code{algo new NAME} & Creates \file{NAME.algo} from a ready-to-run three-agent template\\
\code{algo doctor} & Checks Python, the parser, live-mode readiness and a compiler self-test\\
\code{algo notebook} & Opens \file{AlgoLang\_Notebook.ipynb} in Jupyter, installing Jupyter the first time\\
\code{algo} & Numbered menu of every option\\
\code{algo shell} & Interactive shell: type declarations one at a time\\
\code{algo studio [--port N]} & AlgoLang Studio, the graphical IDE in your browser\\
\code{algo ask "…" [-o F] [--run]} & AI writes a program from a description (compile-and-repair loop)\\
\code{algo wizard} & Build a program by answering questions, no AI needed\\
\code{algo --version} & Prints the version\\
\bottomrule
\end{tabularx}
\end{center}

\section{Options for \texttt{run}}

\begin{center}
\small
\begin{tabularx}{\linewidth}{@{}L{4.4cm}Y@{}}
\toprule
\thead{Option} & \thead{Meaning}\\
\goldrule
\code{--arg "name=value"} & Supplies a flow input; repeat for several inputs. Inputs not supplied are asked for interactively\\
\code{--live} & Runs agents on real Claude models (needs \code{ANTHROPIC\_API\_KEY})\\
\code{--interactive} & Pauses on low-confidence results for human approval\\
\code{--audit-dir DIR} & Writes audit logs to \code{DIR} (default \file{runs})\\
\bottomrule
\end{tabularx}
\end{center}

\begin{terminal}
$ algo run examples/paper_review.algo \
      --arg "paper=LG-CEF: lesion-guided ensemble for DR grading" \
      --interactive --audit-dir audit/2026-10
\end{terminal}

\section{Exit codes}

\begin{center}
\small
\begin{tabularx}{\linewidth}{@{}L{1.3cm}Y@{}}
\toprule
\thead{Code} & \thead{Meaning}\\
\goldrule
0 & Success\\
1 & Compile errors, a missing file or input, or a malformed option\\
2 & A policy violation stopped the run\\
3 & An agent call failed after its retries\\
\bottomrule
\end{tabularx}
\end{center}

Scripts and continuous-integration jobs can rely on these codes; for example, a CI job can run \code{algo check} on every \file{.algo} file and fail the build on exit code 1.

\section{Environment variables}

\begin{center}
\small
\begin{tabularx}{\linewidth}{@{}L{4.3cm}Y@{}}
\toprule
\thead{Variable} & \thead{Meaning}\\
\goldrule
\code{ANTHROPIC\_API\_KEY} & API key for live mode\\
\code{ALGO\_HOME} & Folder containing the runtime; set automatically by \code{algo}\\
\code{ALGO\_LIVE=1} & Live mode (set by \code{--live})\\
\code{ALGO\_INTERACTIVE=1} & Interactive human review (set by \code{--interactive})\\
\code{ALGO\_AUDIT\_DIR} & Audit-log folder (set by \code{--audit-dir})\\
\bottomrule
\end{tabularx}
\end{center}

\section{Running compiled programs directly}

\code{algo build} produces an ordinary Python script that can be run, scheduled or shipped without the compiler:

\begin{terminal}
$ algo build research_brief.algo
✓ compiled research_brief.algo → build/research_brief.py
$ python build/research_brief.py --topic "Federated learning in IoT"
\end{terminal}

The script needs only \file{algoruntime.py}, found through \code{ALGO\_HOME} or the path recorded at compile time.
